\documentclass[10pt,twoside]{article}

\usepackage{amssymb}
\usepackage{amsmath}
\usepackage{latexsym}
\usepackage{color}
\usepackage{graphics}
\usepackage{xspace}

% Commands for special characters
\newcommand{\mybackslash}{\char'134}
\newcommand{\myleftbracket}{\char'173}
\newcommand{\myrightbracket}{\char'175}
\newcommand{\mypercent}{\char'045}
\newcommand{\myunderscore}{\char'137}

\usepackage{graphicx}
\usepackage{libertine}

% The 'ifthen' package supports Boolean flags
\usepackage{ifthen}
% The 'solutions' flag determines whether this is the original handout
%    or a solution
\newboolean{solutions}
\setboolean{solutions}{False}  % Default is original handout

% Uncomment the next line before starting on the solutions
% \setboolean{solutions}{True}

\newcommand{\coursenumber}{CS 5594}
\newcommand{\coursetitle}{Blockchain Technologies}
\newcommand{\docdate}{Feb 09, 2024}
\newcommand{\duedate}{Feb 22, 2024}
\newcommand{\homeworknumber}{1}

% The document title depends on whether these are solutions
\ifthenelse{\boolean{solutions}}{% Then
	\newcommand{\doctitle}{Solutions to Homework Assignment \homeworknumber}
}{% Else
	\newcommand{\doctitle}{Homework Assignment \homeworknumber}
}

% The document date depends on whether these are solutions
\ifthenelse{\boolean{solutions}}{% Then
	\renewcommand{\docdate}{\duedate}
}{% Else
}

% If you are the author, so put your name here
\renewcommand{\author}{Thang Hoang}

\pagestyle{myheadings}
\markboth{\hfill\doctitle\hfill\docdate}{\docdate\hfill\doctitle\hfill}

\addtolength{\textwidth}{1.00in}
\addtolength{\textheight}{1.00in}
\addtolength{\evensidemargin}{-1.00in}
\addtolength{\oddsidemargin}{-0.00in}
\addtolength{\topmargin}{-.50in}
\setlength{\footskip}{0pt}

\newcommand{\polyreduce}{\leq_{\mathrm{P}}}

\newcounter{problem}
\setcounter{problem}{0}
\newcommand{\problem}[1]{%
	\refstepcounter{problem}\noindent\textbf{[#1] \arabic{problem}.}}

\newcommand{\solution}{\bigskip\hrule\bigskip}
\newcommand{\problembreak}{\bigskip\hrule\bigskip}

\renewcommand{\theenumi}{\textbf{\Alph{enumi}}}
\renewcommand{\theenumii}{\textbf{\roman{enumii}}}

\newcommand{\emptystring}{\lambda}

% Pseudocode comment symbol
\newcommand{\comment}{\textbf{//}\ \ }

% Pseudocode line numbering
\newcounter{pseudocode}
\newcommand{\firstline}{\setcounter{pseudocode}{0}\linenumber}
\newcommand{\linenumber}{\refstepcounter{pseudocode}\thepseudocode}
\newcommand{\pseudoindent}{\hspace*{26pt}}

\newcommand{\nil}{\mbox{\textsc{nil}}}

% Mathematical symbols
\newcommand{\grid}{\mathcal{G}}  % Grid graph
\newcommand{\integers}{\mathbb{Z}}  % Integers
\newcommand{\naturals}{\mathbb{N}}  % Natural numbers
\newcommand{\reals}{\mathbb{R}}  % Real numbers

% Probability
\newcommand{\expect}[1]{\mathbf{E}\left[#1\right]}
\newcommand{\prob}[1]{\mathrm{Pr}\left[#1\right]}
\newcommand{\var}[1]{\mathrm{Var}\left[#1\right]}

% Logic
\newcommand{\NOT}[1]{\neg #1}
\newcommand{\AND}{\wedge}
\newcommand{\OR}{\vee}
\newcommand{\clause}[1]{\left(#1\right)}

\newlength{\problemoffset}
\setlength{\problemoffset}{0.5in}

% Decision problem macro
% A command for formatting decision problems a la Garey and Johnson
\newcommand{\decision}[3]{%     \decision{NAME}{INSTANCE}{QUESTION}
	\begin{list}{}{
			\setlength{\leftmargin}{\problemoffset}
			\setlength{\rightmargin}{\problemoffset}
			\setlength{\parsep}{0pt}
			\setlength{\itemsep}{2pt}
			\setlength{\topsep}{\itemsep}
			\setlength{\partopsep}{\itemsep}
		}
		\item
		{\textsc{#1}}
		\item
		{INSTANCE: #2}
		\item
		{QUESTION: #3}
	\end{list}
}

% Optimization problem macro
\newcommand{\optimization}[3]{%  \optimization{NAME}{INSTANCE}{SOLUTION}
	\begin{list}{}{
			\setlength{\leftmargin}{\problemoffset}
			\setlength{\rightmargin}{\problemoffset}
			\setlength{\parsep}{0pt}
			\setlength{\itemsep}{2pt}
			\setlength{\topsep}{\itemsep}
			\setlength{\partopsep}{\itemsep}
		}
		\item
		{\rule{0pt}{14pt}\textsc{#1}}
		\item
		{INSTANCE: #2}
		\item
		{SOLUTION: #3}
	\end{list}
}

\newcommand{\reaches}{\leadsto}

% Table layout
\newcommand{\toprule}{\rule[11pt]{0pt}{2pt}}
\newcommand{\bottomrule}{\rule[-6pt]{0pt}{0pt}}
\newenvironment{raggedpars}[1][2.0in]{%
	\begin{minipage}[t]{#1}\raggedright\toprule}%
	{\bottomrule\end{minipage}}

% Dynamic programming macros
\newlength{\arrowwidth}
\setbox3=\hbox{$\nwarrow$}
\setlength{\arrowwidth}{\wd3}
\newcommand{\optnwarrow}[1]{\if1#1\nwarrow%
	\else\rule{\arrowwidth}{0pt}\fi}
\newcommand{\optuparrow}[1]{\if1#1\uparrow%
	\else\rule{\arrowwidth}{0pt}\fi}
\newcommand{\optleftarrow}[1]{\if1#1\leftarrow%
	\else\rule{\arrowwidth}{0pt}\fi}
% Use \tablebox to specify any combination of arrows, plus the value
\newcommand{\tablebox}[4]{%
	\setlength{\arraycolsep}{0.0pt}%
	\begin{array}{cc}%
		\optnwarrow{#1} & \optuparrow{#2} \\%
		\optleftarrow{#3} & #4%
	\end{array}%
}
% Use \tableboxred to specify any combination of arrows, plus the value
% The value will be red; arrows are made red if 2 used instead of 1
\newcommand{\optnwarrowred}[1]{\if1#1\nwarrow%
	\else\if2#1\textcolor{red}{\nwarrow}\else\rule{\arrowwidth}{0pt}\fi\fi}
\newcommand{\optuparrowred}[1]{\if1#1\uparrow%
	\else\if2#1\textcolor{red}{\uparrow}\else\rule{\arrowwidth}{0pt}\fi\fi}
\newcommand{\optleftarrowred}[1]{\if1#1\leftarrow%
	\else\if2#1\textcolor{red}{\leftarrow}\else\rule{\arrowwidth}{0pt}\fi\fi}
\newcommand{\tableboxred}[4]{%
	\setlength{\arraycolsep}{0.0pt}%
	\begin{array}{cc}%
		\optnwarrowred{#1} &%
		\optuparrowred{#2} \\%
		\optleftarrowred{#3} &%
		\textcolor{red}{#4}%
	\end{array}%
}

% Allow really sloppy paragraphs that do not generate overfull or
%    underfull hbox's
\newenvironment{SLOPPY}{\begin{sloppypar}\hbadness=10000}{\end{sloppypar}}

% Definitions for this homework
\newcommand{\extent}[1]{\mathrm{extent}(#1)}
\newcommand{\opt}[2]{\mbox{\textsc{opt}}(#1,#2)}
\newcommand{\gap}{\mbox{\texttt{-}}}
\newcommand{\rewriterule}[2]{#1\to #2}
\newcommand{\rewrites}[2][]{\mathop{\Longrightarrow}\limits_{#1}^{#2}}
\newcommand{\reduces}{\leq}
\newcommand{\polyreduces}{\leq_P}
\newcommand{\mathsc}[1]{\mbox{\normalfont\textsc{#1}}}
\newcommand{\NP}{\mathcal{NP}}
\renewcommand{\P}{\mathcal{P}}

\begin{document}
	
	\thispagestyle{empty}
	
	\begin{center}
		\begin{tabular}{lcr}
			\multicolumn{3}{c}{\Large\textbf{\coursenumber}}
			\\[0.04in]
			\multicolumn{3}{c}{\Large\textbf{\doctitle}}
			% If these are solutions, then include the author's (student's) name
			\ifthenelse{\boolean{solutions}}{% Then
				\\[0.04in]\multicolumn{3}{c}{\large\textbf{\author}}
			}{} % Else omit
			% If these are solutions, then the date is the due date
			\ifthenelse{\boolean{solutions}}{% Then
				\\[0.10in]\multicolumn{3}{c}{\duedate}
			}{% Else, put given and due dates
				\\[0.10in]
				\textbf{Given:} \docdate
				& \hspace*{1.0in} &
				\textbf{Due:} \duedate
			}
		\end{tabular}
	\end{center}
	
	% If these are solutions, then we do not include the directions
	\ifthenelse{\boolean{solutions}}{} % No directions
	{
		% Original document includes directions
		\begingroup % This allows an argument that contains multiple paragraphs
		\paragraph{General directions.}
		The point value of each problem is shown in [ ].
		Each solution must include all details and
		an explanation of why the given solution is correct.
		\textbf{\textcolor{red}{In particular,
				write complete sentences.
				A correct answer without an explanation is worth no credit.}}
		The completed assignment must be submitted on Canvas as a PDF by 11:59 PM
		on \duedate.
		\textbf{No late homework will be accepted.}
		
		\paragraph{Digital preparation of your solutions is mandatory.}
		Use of \LaTeX\ is optional, but encouraged.
		No matter how you prepare your homework,
		\textbf{please include your name.}
		
		
		\problembreak
		
		\newpage
		
	}
	
	\problem{20} Consider the three attributes including Consistency, Availability, and Partition Tolerance in a distributed system:
	{\bfseries 
		\begin{enumerate}
			
			\item Prove that a distributed system cannot achieve these three attributes simultaneously.
			
				
			\item Show how public blockchain manages to achieve these attributes. Specifically, which attribute is scarified in favor of the other two attributes? Describe the core technique that the public blockchain used to fully achieve those two attributes, while managing to (eventually) offer the remaining one. 
					
			
		\end{enumerate}
	}
	
	\solution
	
	% PUT YOUR SOLUTION HERE
	
	\problembreak
	
	\problem{50} Suppose a sender $ S $ uses the following Merkle-hash tree $ T $ to authenticate messages $(M_1,\dots,M_8)$ to a receiver $ R $.
		
		\begin{figure}[!h]
			\centering
			\includegraphics[width=.8\textwidth]{merke-tree.pdf}
		\end{figure}
		
		{\bfseries
			\begin{enumerate}
				\item How would one authenticate message $ M_4 $? What elements of $ T $ must be transmitted from $ S $ to $ R $, and write the correct verification equation.
				
				\item Why the Merkle-hash tree $ T $ has an additional level of hash in the leaves?  
				
				\item What are two necessary conditions for a set of data to be authenticated by the Merkle-hash tree? 
				
				\item Show how to find a collision in a Merkle-hash tree $T'$ with a \emph{flexible} structure (i.e., the number of the inputs is not fixed). Specifically show how to find two sets of messages $ \mathcal{A}=\{A_1,\dots,A_t\}$ and $\mathcal{B} = \{B_1,\dots,B_{2t}\} $ such that $\mathsf{MerkleRoot}(\mathcal{A}) = \mathsf{MerkleRoot}(\mathcal{B})$.
				
				\item Describe how Merkle-hash trees are used to achieve integrity in public blockchain (e.g., bitcoin). 
				
			\end{enumerate}
		}
				
	\solution

	% PUT YOUR SOLUTION HERE
	
	\problembreak

	
	
	\problem{30} In class, we have covered a Digital Signature Algorithm (DSA), in which the signature $(r,s)$ for the message $m$ can be computed as
	
	\begin{equation*}
		\begin{cases}
			r=(g^k \mod p) \mod q\\
			s = k^{-1}  (H(m) + x \cdot r) \mod q
		\end{cases}
	\end{equation*}
	where $k$ is a random private key per signing, $x$ is long-term private key and $H$ is a cryptographic hash function.
	{\bfseries
	\begin{enumerate}
		\item Consider a variant of DSA algorithm, in which the second component of the signature generation is computed as 
		
		$$ s = k^{-1}  (m + x \cdot r) \mod q$$
		Show that this variant is not secure, in which the attacker can forge valid signature for any arbitrary message of it choice without querying any signatures from the signer. 
		
		\item 	Sony PS3 was hacked by the hacker group ``\textsf{fail0Overflow}'' via a key recovery attack on the ECDSA digital signatures computed in Sony PS3 platform.
		%
		Explain what caused the attack and show the steps of the attack in details.  
	\end{enumerate}
	}
		
		\solution
	
		% PUT YOUR SOLUTION HERE
		
		\problembreak
		
	
		\problem{25} Consider the following ECC curve $E$:
	$$ Y^2 =X^3 +231X+473, p=17389, q=1321, G=(11259,11278) \in E(\mathbb{F}_p).
	$$ 
	{\bfseries
		\begin{enumerate}
			\item Assume the signing key of Alice is $ sk = 542 $. What is her corresponding public key $pk$? What is her signature on the hash of a message $ H(m) = 644 $ with the ephemeral key $ k = 847 $? 
			
			\item Assume the public key of Bob is $ pk = (14594,308) $. What is his private key $ sk $? (you can use any method to find it, but describe it in details). Use his private key $ sk $ that you found to forge his signature on the hash of a message $ H(m) = 516 $ using the ephemeral key $ k = 365 $. 
			
		\end{enumerate}
	}
	\solution

	% PUT YOUR SOLUTION HERE
	
	\problembreak
\end{document}

